\section{Introduction}

Vision transformers develop a peculiar failure of interpretability at scale.
Darcet \textit{et al.}~\cite{darcet2024registers} observed that in large
supervised and self-supervised ViTs a small number of patch tokens acquire
anomalously large activation norms, and that the \texttt{[CLS]} token routes a
disproportionate share of its attention to them. These tokens sit in
low-information image regions, carry global rather than local content, and make
attention maps unusable as saliency evidence. The proposed remedy is
architectural and cheap: prepend $K$ learnable \emph{register} tokens to the
input sequence, discard them at the output, and let the model use them as the
scratch space it was previously carving out of the image itself. At scale this
works---the artifacts disappear and downstream performance is unaffected or
mildly improved.

The mechanism reported in that work is a \emph{cleanup} effect, demonstrated in
a regime where the model has far more data than it can overfit. This leaves an
obvious question unanswered. If registers absorb capacity that the model would
otherwise spend on memorising the training set, they should behave like a
regularizer, and the place to look for a regularizer is the regime where
overfitting is the binding constraint: a small model trained on scarce data.

We therefore ask a narrow, falsifiable question:

\begin{quote}
\textit{When a ViT is starved of data, do register tokens reduce the
generalization gap, and does any such reduction transfer to held-out
performance?}
\end{quote}

We test three hypotheses drawn from the original mechanism:

\begin{itemize}
\item[\textbf{H1}] \textbf{Structural regularization.} Registers reduce the
generalization gap $\Delta\mathcal{L} = \mathcal{L}_{\mathrm{val}} -
\mathcal{L}_{\mathrm{train}}$.
\item[\textbf{H2}] \textbf{Capacity dilution.} At an embedding width of
$d = 192$, a large $K$ consumes representational capacity and degrades
accuracy, producing a non-monotonic response in $K$.
\item[\textbf{H3}] \textbf{Entropy collapse prevention.} Registers prevent the
layer-wise collapse of attention entropy that accompanies artifact formation.
\end{itemize}

Our contributions are:

\begin{enumerate}
\item A fully paired $4 \times 3$ ablation ($K \in \{0,1,4,8\}$, seeds
\texttt{42}, \texttt{1337}, \texttt{3407}) on a strict 100-images-per-class
CIFAR-100 subset, with every hyperparameter, data split and schedule held fixed
across arms, so that $K$ is the only manipulated variable.
\item A statistical treatment appropriate to $n = 3$: differences are formed
\emph{within} a seed before reduction, which removes the seed variance that
otherwise dominates the treatment effect at this sample size.
\item A separation of validation-measured from test-measured effects. This
separation is what turns an apparently positive result into a negative one: the
gap reduction predicted by H1 is present and consistent on the validation
split, and absent on the held-out test split.
\item A negative result reported as such, with the mechanism-level evidence
(layer-wise entropy, patch-norm outlier rates, training-loss equality) that
distinguishes ``registers did nothing'' from ``registers helped somewhere we
could not measure''.
\end{enumerate}

We find no configuration that beats the register-free control on held-out
accuracy. H1 is not supported on held-out data, H2 is supported at $K=8$, and
H3 is not supported. Section~\ref{sec:discussion} argues that the most likely
reason is a regime mismatch rather than a refutation of
\cite{darcet2024registers}: our backbone is ImageNet-pretrained, the registers
are randomly initialised into it, and 50 epochs over 9{,}000 images is a short
window in which to learn a new routing behaviour from scratch.
